\documentclass[12pt,a4paper]{article}

% Kódování a český jazyk
\usepackage[T1]{fontenc}
\usepackage[utf8]{inputenc}
\usepackage[czech]{babel}
\usepackage{lmodern}

% Vzhled a sazba
\usepackage[a4paper,margin=2.5cm]{geometry}
\usepackage{setspace}
\usepackage{booktabs}
\usepackage{array}
\usepackage{hyperref}
\usepackage{xurl}

\hypersetup{
    colorlinks=true,
    linkcolor=black,
    urlcolor=blue,
    citecolor=black,
    pdftitle={Bezpečnost hesel a ochrana uživatelských účtů},
    pdfauthor={Doplňte jméno}
}

\onehalfspacing
\setlength{\parindent}{1.25cm}
\setlength{\parskip}{0.25em}

% ===== ÚDAJE K DOPLNĚNÍ =====
\newcommand{\nazevprace}{Bezpečnost hesel a ochrana uživatelských účtů}
\newcommand{\jmeno}{Doplňte jméno a příjmení}
\newcommand{\predmet}{Doplňte název předmětu}
\newcommand{\skola}{Doplňte název školy}
\newcommand{\datumodevzdani}{\today}

\begin{document}

\begin{titlepage}
    \centering
    {\large \skola\par}
    \vspace{3cm}
    {\LARGE\bfseries \nazevprace\par}
    \vspace{2cm}
    {\large Praktické cvičení: Práce s více soubory v LaTeXu\par}
    \vfill
    \begin{flushleft}
        \textbf{Autor/ka:} \jmeno\\[0.5em]
        \textbf{Předmět:} \predmet\\[0.5em]
        \textbf{Datum:} \datumodevzdani
    \end{flushleft}
    \vspace{1cm}
\end{titlepage}

\tableofcontents
\newpage

\section{Úvod}

% TUTO ČÁST NAPIŠ VLASTNÍMI SLOVY:
% Vysvětli, proč lidé používají hesla, jaké účty chrání a proč je jejich
% zabezpečení důležité. Uveď, čemu se bude práce věnovat.
% Vhodné je zde odkázat na odborné doporučení, například \cite{nist2025}.

\section{Zásady tvorby bezpečných hesel}

% TUTO ČÁST NAPIŠ VLASTNÍMI SLOVY:
% Popiš, proč je důležitá délka hesla, jeho jedinečnost a proč není vhodné
% používat stejné heslo pro více služeb. Aktuální doporučení NIST uvádí
% pro hesla používaná jako jediný faktor minimální délku 15 znaků
% \cite{nist2025}. Formuluj závěry vlastními slovy.

\subsection{Délka a jedinečnost hesla}

% Doplň vlastní odstavec o dlouhých heslových frázích, běžných heslech
% a riziku opakovaného použití hesla. Nezaměňuj délku za pouhé střídání
% velkých písmen, číslic a symbolů.

\subsection{Správci hesel a vícefaktorové ověřování}

% Doplň vlastní odstavec o tom, jak správce hesel pomáhá vytvářet a uchovávat
% jedinečná hesla. Vysvětli také princip vícefaktorového ověřování (MFA).
% K doporučením můžeš využít NIST a technické pokyny ENISA \cite{enisa2025}.

\section{Nejčastější rizika a možnosti ochrany}

% TUTO ČÁST NAPIŠ VLASTNÍMI SLOVY:
% Vysvětli například phishing, únik přihlašovacích údajů a útoky,
% při kterých útočníci zkoušejí uniklá hesla u dalších služeb.
% Poté ručně sestav seznam hlavních preventivních zásad podle zdrojů.

% Seznam, který máš doplnit vlastními slovy:
% \begin{itemize}
%   \item ...
%   \item ...
%   \item ...
% \end{itemize}

% Následující tabulku můžeš použít jako AI připravenou technickou část.
\begin{table}[ht]
    \centering
    \caption{Přehled vybraných způsobů ochrany uživatelských účtů}
    \label{tab:ochrana}
    \begin{tabular}{p{3.2cm} p{5.0cm} p{5.0cm}}
        \toprule
        \textbf{Opatření} & \textbf{Přínos} & \textbf{Na co si dát pozor} \\
        \midrule
        Jedinečné heslo pro každou službu
        & Omezuje dopad úniku hesla z jedné služby na ostatní účty.
        & Hesla je obtížné bezpečně spravovat bez vhodného nástroje. \\
        \addlinespace
        Správce hesel
        & Pomáhá vytvářet a uchovávat dlouhá, náhodná a odlišná hesla.
        & Hlavní přístup k trezoru musí být dobře zabezpečen. \\
        \addlinespace
        Vícefaktorové ověřování (MFA)
        & Přidává další ověřovací prvek nad rámec samotného hesla.
        & Ne všechny metody poskytují stejnou úroveň ochrany; vhodné je preferovat odolné metody. \\
        \addlinespace
        Opatrnost vůči phishingu
        & Snižuje riziko předání přihlašovacích údajů podvodné stránce.
        & Je nutné ověřovat adresu webu a neočekávané výzvy k přihlášení. \\
        \bottomrule
    \end{tabular}
\end{table}

\section{Závěr}

% TUTO ČÁST NAPIŠ VLASTNÍMI SLOVY:
% Shrň nejdůležitější poznatky. Uveď, která opatření považuješ za
% nejpraktičtější a proč. Nepřidávej zde nové zdroje ani nové téma.


\newpage
\begin{thebibliography}{9}

\bibitem{nist2025}
TEMOSHOК, David, et al.
\textit{Digital Identity Guidelines: Authentication and Authenticator Management}.
NIST Special Publication 800-63B-4. Gaithersburg: National Institute of Standards and Technology, 2025.
Dostupné z: \url{https://pages.nist.gov/800-63-4/sp800-63b.html}

\bibitem{enisa2025}
EUROPEAN UNION AGENCY FOR CYBERSECURITY (ENISA).
\textit{Technical Implementation Guidance on Cybersecurity Risk Management Measures}.
Version 1.0. 2025.
Dostupné z: \url{https://www.enisa.europa.eu/sites/default/files/2025-06/ENISA_Technical_implementation_guidance_on_cybersecurity_risk_management_measures_version_1.0.pdf}

\bibitem{anderson2020}
ANDERSON, Ross.
\textit{Security Engineering: A Guide to Building Dependable Distributed Systems}.
3rd ed. Hoboken: Wiley, 2020.

\end{thebibliography}


\end{document}
